\section{The native full-flux magnetic sector}
\label{cusp-native-full-flux-section}

This section retains the canonical radial parameter, period matrix, Euclidean
quotient metric, marked lattice, and every sign already fixed in
\cref{cusp-spec-section}.  It treats the monodromy-invariant integral class
that uses both elementary alternating blocks.  It is therefore distinct from
the unit-flux plane treated in \cref{cusp-magnetic-section}.  No metric
rescaling, finite covering, rectangular replacement of the lattice, or
orientation suppression occurs below.

For \(T>T_0\), retain
\begin{equation}
 \begin{gathered}
 L_T=T+H_T=\operatorname{Im}\tau_T,
 \qquad q_T=L_T-B_T=-\operatorname{Im}\beta_T,\\
 m_T=\operatorname{Im}\mu_T,
 \qquad D_T=L_Tq_T+6m_T^2>0,
 \end{gathered}
 \label{cusp-native-full-flux-canonical-scalars}
\end{equation}
and
write
\[
 \eta_T^{\mathrm{re}}:=\operatorname{Re}h_T.
\]
This scalar is denoted by \(\eta_T\) in
\cref{cusp-spec-section}; the superscript is retained throughout this
section so that it cannot be confused with the native integral two-form
introduced below.  The complete period matrix is
\begin{equation}
 P_T=
 \begin{pmatrix}
  6\alpha_T&\eta_T^{\mathrm{re}}&1&0\\
  6m_T&L_T&0&0\\
  c_T&\alpha_T&0&1\\
  -q_T&m_T&0&0
 \end{pmatrix},
 \qquad \det P_T=-D_T.
 \label{cusp-native-full-flux-period-matrix}
\end{equation}
Thus \(F_T=\mathbb R^4/P_T\mathbb Z^4\) has the unrescaled Euclidean quotient
metric and volume \(D_T\).  In the marking coordinates \(y\in\mathbb R^4\),
the map \(y\mapsto x=P_Ty\) pulls the Euclidean metric back to
\(G_T=P_T^{\mathsf T}P_T\).

The limiting cusp constants remain
\begin{equation}
 m=\operatorname{Im}\mu(0),\qquad
 H=\operatorname{Im}h(0),\qquad
 B=\operatorname{Im}b(0),
 \label{cusp-native-full-flux-limiting-constants}
\end{equation}
so holomorphy at \(e^{-2\pi T}=0\) gives
\begin{equation}
 \begin{aligned}
 m_T&=m+O(e^{-2\pi T}),&
 H_T&=H+O(e^{-2\pi T}),\\
 B_T&=B+O(e^{-2\pi T}),&
 D_T&=T^2+(2H-B)T\\
 &&&\quad +(H^2-BH+6m^2)+O(Te^{-2\pi T}).
 \end{aligned}
 \label{cusp-native-full-flux-cusp-expansions}
\end{equation}

\subsection{The invariant class, its cocycle, and both orientations}

In the ordered marking
\(\mathbb Z\langle\gamma,u,w,\delta\rangle\), put
\begin{equation}
 \eta=u\wedge w+6\gamma\wedge\delta.
 \label{cusp-native-full-flux-class}
\end{equation}
The two finite monodromy generators act by
\begin{align*}
T_1\gamma&=\gamma,&T_1u&=-u-w,&
T_1w&=-6\gamma+u,&T_1\delta&=2\gamma+u+w+\delta,\\
T_2\gamma&=\gamma,&T_2u&=6\gamma+w,&
T_2w&=-u,&T_2\delta&=-3\gamma+u+\delta.
\end{align*}
Directly, without passing to a normalized representative,
\begin{align*}
(\bigwedge^2T_1)\eta
&=(-u-w)\wedge(-6\gamma+u)
 +6\gamma\wedge(2\gamma+u+w+\delta)\\
&=u\wedge w+6\gamma\wedge\delta,\\
(\bigwedge^2T_2)\eta
&=(6\gamma+w)\wedge(-u)
 +6\gamma\wedge(-3\gamma+u+\delta)\\
&=u\wedge w+6\gamma\wedge\delta.
\end{align*}
Hence \(\eta\) is invariant under both generators.

On \(Y=\mathbb R^4/\mathbb Z^4\), with ordered coordinates
\(y^1,y^2,y^3,y^4\), its constant representative and alternating matrix are
\begin{equation}
 \eta_T=\mathrm dy^2\wedge\mathrm dy^3
       +6\,\mathrm dy^1\wedge\mathrm dy^4,
 \qquad
 E_\eta=
 \begin{pmatrix}
 0&0&0&6\\
 0&0&1&0\\
 0&-1&0&0\\
 -6&0&0&0
 \end{pmatrix}.
 \label{cusp-native-full-flux-form-matrix}
\end{equation}
In the paired basis \((e_2,e_3;e_1,e_4)\), the alternating elementary
divisors are \((1,6)\); consequently
\begin{equation}
 \operatorname{SNF}(E_\eta)=\operatorname{diag}(1,1,6,6),
 \qquad \operatorname{Pf}(E_\eta)=6,
 \qquad \det E_\eta=36.
 \label{cusp-native-full-flux-smith-pfaffian}
\end{equation}

The actual Hermitian line bundle is the quotient of
\(\mathbb R^4\times\mathbb C\) by
\begin{align}
 (y,z)&\longmapsto(y+e_1,e^{12\pi i y^4}z),&
 (y,z)&\longmapsto(y+e_2,e^{2\pi i y^3}z),
 \label{cusp-native-full-flux-cocycle}\\
 (y,z)&\longmapsto(y+e_3,z),&
 (y,z)&\longmapsto(y+e_4,z).
 \notag
\end{align}
The only possibly nontrivial generator commutators acquire the factors
\(e^{12\pi i}=1\) and \(e^{2\pi i}=1\), so these maps define a
\(\mathbb Z^4\)-action.  Denote the quotient by
\(\mathscr L_{\eta,T}\to F_T\), using the marking
\([y]\mapsto[P_Ty]\).  The connection
\begin{equation}
 \nabla_{\eta,T}
 =\mathrm d-2\pi i(6y^1\,\mathrm dy^4+y^2\,\mathrm dy^3)
 \label{cusp-native-full-flux-connection}
\end{equation}
descends through the displayed cocycle.  For instance translation by \(e_1\)
changes its connection form by \(-12\pi i\,\mathrm dy^4\), which equals the
negative logarithmic derivative of \(e^{12\pi i y^4}\); the \(e_2\) check is
identical.  Therefore
\begin{equation}
 F_{\nabla_{\eta,T}}=-2\pi i\eta_T,
 \qquad c_1(\mathscr L_{\eta,T})
 =\frac{i}{2\pi}F_{\nabla_{\eta,T}}=\eta_T.
 \label{cusp-native-full-flux-curvature}
\end{equation}

There are two opposite orientations which must not be conflated.  Write
\begin{equation}
 \operatorname{or}_{\mathrm{mark}}
 =\mathrm dy^1\wedge\mathrm dy^2\wedge
  \mathrm dy^3\wedge\mathrm dy^4
 \label{cusp-native-full-flux-marking-orientation}
\end{equation}
and let \(\operatorname{or}_{x}\) be the native complex orientation of
\((z_1,z_2)=(x_1+iy_1,x_2+iy_2)\).  Then
\begin{equation}
 P_T^*\operatorname{or}_{x}
 =-D_T\operatorname{or}_{\mathrm{mark}},
 \label{cusp-native-full-flux-orientation-reversal}
\end{equation}
whereas
\begin{equation}
 \frac12\eta_T^2=6\operatorname{or}_{\mathrm{mark}}.
 \label{cusp-native-full-flux-square}
\end{equation}
Consequently
\begin{equation}
 \frac12\langle\eta_T^2,[F_T]_{\mathrm{mark}}\rangle=6,
 \qquad
 \frac12\langle\eta_T^2,[F_T]_{x\text{-}\mathrm{complex}}\rangle=-6.
 \label{cusp-native-full-flux-two-orientation-pairings}
\end{equation}

\subsection{The exact physical frequencies}

In the physical Euclidean coordinates the two-form matrix is
\begin{equation}
 \Omega_T=P_T^{-\mathsf T}E_\eta P_T^{-1}
 =\frac1{D_T}
 \begin{pmatrix}
 0&-q_T&0&-6m_T\\
 q_T&0&6m_T&0\\
 0&-6m_T&0&6L_T\\
 6m_T&0&-6L_T&0
 \end{pmatrix}.
 \label{cusp-native-full-flux-physical-matrix}
\end{equation}
After the exact permutation
\((x^1,x^2,x^3,x^4)\mapsto(x^1,x^3;x^2,x^4)\), its off-diagonal block is
\begin{equation}
 K_T=\begin{pmatrix}-q_T&-6m_T\\-6m_T&6L_T\end{pmatrix},
 \qquad \det K_T=-6D_T.
 \label{cusp-native-full-flux-K}
\end{equation}
Define the magnetic scalars
\begin{equation}
 S_T:=6L_T-q_T,
 \qquad
 \Xi_T:=\sqrt{(6L_T+q_T)^2+144m_T^2}.
 \label{cusp-native-full-flux-S-Xi}
\end{equation}
Here \(S_T\) has no argument and is distinct from the affine residual
\(S_T(k)\) used in the dual-lattice calculation of
\cref{cusp-spec-exact-dual-lattice}.  The signed eigenvalues of \(K_T\) are
\begin{equation}
 \kappa_{T,\pm}=\frac{S_T\pm\Xi_T}{2}.
 \label{cusp-native-full-flux-kappas}
\end{equation}
Direct expansion, with no asymptotic replacement, gives
\begin{equation}
 \Xi_T^2-S_T^2
 =24(L_Tq_T+6m_T^2)=24D_T>0.
 \label{cusp-native-full-flux-Xi-identity}
\end{equation}
Thus \(\kappa_{T,+}>0>\kappa_{T,-}\).  The two positive charge-one
magnetic frequencies are
\begin{equation}
 \boxed{
 \mathcal B_{T,+}=\frac{\pi(S_T+\Xi_T)}{D_T},
 \qquad
 \mathcal B_{T,-}=\frac{\pi(\Xi_T-S_T)}{D_T}.}
 \label{cusp-native-full-flux-frequencies}
\end{equation}
They obey the exact identities
\begin{align}
 \mathcal B_{T,+}\mathcal B_{T,-}
 &=\frac{24\pi^2}{D_T},
 \label{cusp-native-full-flux-frequency-product}\\
 \mathcal B_{T,+}+\mathcal B_{T,-}
 &=\frac{2\pi\Xi_T}{D_T},&
 \mathcal B_{T,+}-\mathcal B_{T,-}
 &=\frac{2\pi S_T}{D_T},
 \label{cusp-native-full-flux-frequency-sum-difference}\\
 \mathcal B_{T,+}&=\frac{24\pi}{\Xi_T-S_T},&
 \mathcal B_{T,-}&=\frac{24\pi}{\Xi_T+S_T}.
 \label{cusp-native-full-flux-rationalized-frequencies}
\end{align}
The last line follows by multiplying by the conjugate factor and using
\cref{cusp-native-full-flux-Xi-identity}.

Keeping the three independent limiting coefficients \(H,B,m\), one obtains
\begin{align}
 \mathcal B_{T,-}
 &=\frac{2\pi}{T}-\frac{2\pi H}{T^2}
 +\frac{2\pi}{T^3}\left(H^2-\frac67m^2\right)+O(T^{-4}),
 \label{cusp-native-full-flux-Bminus-expansion}\\
 \mathcal B_{T,+}
 &=\frac{12\pi}{T}+\frac{12\pi(B-H)}{T^2}
 +\frac{12\pi}{T^3}
 \left((B-H)^2-\frac{36}{7}m^2\right)+O(T^{-4}),
 \label{cusp-native-full-flux-Bplus-expansion}\\
 \rho_T:=\frac{\mathcal B_{T,+}}{\mathcal B_{T,-}}
 &=6+\frac{6B}{T}
 +\frac{6B^2-6BH-180m^2/7}{T^2}+O(T^{-3}).
 \label{cusp-native-full-flux-ratio-expansion}
\end{align}
These follow by substituting \cref{cusp-native-full-flux-cusp-expansions}
into \cref{cusp-native-full-flux-rationalized-frequencies} and expanding the
reciprocals; no sign of \(H,B,m\) is assumed.  In particular, on a sufficiently
deep cusp tail,
\begin{equation}
 S_T>0,
 \qquad \mathcal B_{T,+}>\mathcal B_{T,-}>0,
 \qquad 5<\rho_T<7.
 \label{cusp-native-full-flux-deep-tail-window}
\end{equation}

\subsection{Joint Landau levels and their exact multiplicities}

Let \(\Delta_{\eta,T}=\nabla_{\eta,T}^*\nabla_{\eta,T}\) be the Friedrichs
magnetic Laplacian on \(L^2(F_T,\mathscr L_{\eta,T})\).  Its form domain is
\(H^1(F_T;\mathscr L_{\eta,T})\), its operator domain is
\(H^2(F_T;\mathscr L_{\eta,T})\), and compactness of the torus gives compact
resolvent.  We now prove the multiplicity rather than infer it from the
Pfaffian alone.

Define \(A_T^{\mathrm{pol}}\) on the marked tangent space by
\begin{equation}
 G_T(A_T^{\mathrm{pol}}v,w)=2\pi E_\eta(v,w),
 \qquad
 Q_T^{\mathrm{pol}}=\bigl(-(A_T^{\mathrm{pol}})^2\bigr)^{1/2},
 \qquad
 J_{\eta,T}=A_T^{\mathrm{pol}}(Q_T^{\mathrm{pol}})^{-1}.
 \label{cusp-native-full-flux-polar-complex-structure}
\end{equation}
The endomorphism \(A_T^{\mathrm{pol}}\) is \(G_T\)-skew-adjoint and invertible,
so \(Q_T^{\mathrm{pol}}\) is positive, commutes with
\(A_T^{\mathrm{pol}}\), and
\begin{equation}
 J_{\eta,T}^2=-1,
 \qquad
 E_\eta(J_{\eta,T}v,J_{\eta,T}w)=E_\eta(v,w),
 \qquad
 2\pi E_\eta(v,J_{\eta,T}v)
 =G_T(Q_T^{\mathrm{pol}}v,v)>0.
 \label{cusp-native-full-flux-polar-properties}
\end{equation}
For every signed charge \(q\in\mathbb Z\setminus\{0\}\), use the convention
\begin{equation}
 \mathscr L_{\eta,T}^{\otimes q}=
 \begin{cases}
  (\mathscr L_{\eta,T})^{\otimes q},&q>0,\\
  (\mathscr L_{\eta,T}^{\vee})^{\otimes(-q)},&q<0,
 \end{cases}
 \qquad
 A_{q,T}^{\mathrm{pol}}=qA_T^{\mathrm{pol}},
 \qquad
 Q_{q,T}^{\mathrm{pol}}
 =\bigl(-(A_{q,T}^{\mathrm{pol}})^2\bigr)^{1/2}
 =|q|Q_T^{\mathrm{pol}},
 \label{cusp-native-full-flux-signed-charge-bundle}
\end{equation}
and define
\begin{equation}
 J_{q,T}=A_{q,T}^{\mathrm{pol}}(Q_{q,T}^{\mathrm{pol}})^{-1}
 =\operatorname{sgn}(q)J_{\eta,T}.
 \label{cusp-native-full-flux-signed-charge-complex-structure}
\end{equation}
All these coefficients are constant, hence the Nijenhuis tensor of
\(J_{q,T}\) vanishes and the marking transports it to a complex structure on
\(F_T\).  The induced connection \(\nabla_{\eta,T}^{(q)}\) has
\begin{equation}
 F_{\nabla_{\eta,T}^{(q)}}=-2\pi i q\eta_T,
 \qquad
 2\pi qE_\eta(v,J_{q,T}v)
 =|q|G_T(Q_T^{\mathrm{pol}}v,v)>0.
 \label{cusp-native-full-flux-signed-charge-positivity}
\end{equation}
Thus its curvature is of type \((1,1)\), and
\((\nabla_{\eta,T}^{(q)})^{0,1}\) makes
\(\mathscr L_{\eta,T}^{\otimes q}\) a positive Appell--Humbert line bundle on
the auxiliary complex torus \((F_T,J_{q,T})\).

The auxiliary orientation is fixed, rather than repaired by an absolute value.
For \(q>0\), compatibility identifies the \(J_{q,T}\)-complex orientation
with the \(E_\eta\)-symplectic orientation; for \(q<0\), replacing \(J\) by
\(-J\) preserves complex orientation in complex dimension two.  Hence
\begin{equation}
 [F_T]_{J_{q,T}}=[F_T]_{\mathrm{mark}}
 =-[F_T]_{x\text{-}\mathrm{complex}}.
 \label{cusp-native-full-flux-auxiliary-native-orientations}
\end{equation}
This auxiliary positive structure is not the native complex structure inherited
from the \(S^6\) construction.  Indeed, with \(E_\eta=\operatorname{Im}H_{\eta,T}\),
the native Appell--Humbert calculation is
\begin{equation}
 G_{\eta,T}=
 \begin{pmatrix}
  2L_T&2m_T\\
  2m_T&-q_T/3
 \end{pmatrix},
 \qquad
 \det G_{\eta,T}=-\frac23D_T<0,
 \qquad H_{\eta,T}=2G_{\eta,T}^{-1}.
 \label{cusp-native-full-flux-native-hermitian-signature}
\end{equation}
Consequently \(H_{\eta,T}\), and also \(qH_{\eta,T}\) for every nonzero
\(q\), has signature \((1,1)\).  The native line bundle is holomorphic because
\(q\eta_T\) is of type \((1,1)\), but it is indefinite and is not the positive
bundle just constructed.

The tangent and canonical bundles of the auxiliary complex torus are trivial.
Kodaira vanishing and Hirzebruch--Riemann--Roch therefore give
\begin{equation}
 \begin{aligned}
 \dim_{\mathbb C}H^0\bigl((F_T,J_{q,T}),
                   \mathscr L_{\eta,T}^{\otimes q}\bigr)
 &=\int_{[F_T]_{J_{q,T}}}\frac{(q\eta_T)^2}{2}
  =\int_{[F_T]_{\mathrm{mark}}}\frac{(q\eta_T)^2}{2}
  =6q^2,\\
 \dim_{\mathbb C}H^0\bigl((F_T,J_{\eta,T}),\mathscr L_{\eta,T}\bigr)
 &=6,\\
 \dim_{\mathbb C}H^0\bigl((F_T,J_{2,T}),
                   \mathscr L_{\eta,T}^{\otimes2}\bigr)
 &=24.
 \end{aligned}
 \label{cusp-native-full-flux-vacuum-dimensions}
\end{equation}
No absolute value occurs in this auxiliary Hirzebruch--Riemann--Roch integral.
Equivalently, the positive polarization has elementary divisors
\((|q|,6|q|)\), and the Appell--Humbert theta theorem gives their product
\(6q^2\); see
\href{https://math.univ-cotedazur.fr/~beauvill/pubs/thetaon.pdf}%
{Beauville, Theorems~2.6 (pp.~107--108) and~3.5 (p.~112)}.
On the native complex torus the same signed Hirzebruch--Riemann--Roch
calculation instead gives
\begin{equation}
 \chi\bigl((F_T,J_{\mathrm{native},T}),
           \mathscr L_{\eta,T}^{\otimes q}\bigr)
 =\int_{[F_T]_{x\text{-}\mathrm{complex}}}
       \frac{(q\eta_T)^2}{2}
 =-6q^2,
 \label{cusp-native-full-flux-native-signed-hrr}
\end{equation}
and positivity vanishing is unavailable there.  Thus the positive auxiliary
section count and the native signed Euler characteristic concern distinct
Dolbeault problems and distinct complex orientations.

Choose a \(G_T\)-orthonormal frame
\(e_+,f_+=J_{\eta,T}e_+,e_-,f_-=J_{\eta,T}e_-\) which diagonalizes
\(Q_T^{\mathrm{pol}}\), labelled so its positive eigenvalues are
\(\mathcal B_{T,+},\mathcal B_{T,-}\).  Put
\begin{equation}
 a_\sigma=
 \frac{\nabla_{e_\sigma}+i\nabla_{f_\sigma}}
      {\sqrt{2\mathcal B_{T,\sigma}}},
 \qquad
 a_\sigma^*=
 \frac{-\nabla_{e_\sigma}+i\nabla_{f_\sigma}}
      {\sqrt{2\mathcal B_{T,\sigma}}},
 \qquad \sigma\in\{+,-\}.
 \label{cusp-native-full-flux-ladders}
\end{equation}
The curvature commutators give
\begin{equation}
 [a_\sigma,a_\tau^*]=\delta_{\sigma\tau},
 \qquad [a_\sigma,a_\tau]=[a_\sigma^*,a_\tau^*]=0,
 \label{cusp-native-full-flux-CCR}
\end{equation}
and, for \(N_\sigma=a_\sigma^*a_\sigma\),
\begin{equation}
 \Delta_{\eta,T}
 =\mathcal B_{T,+}(2N_++1)+\mathcal B_{T,-}(2N_-+1).
 \label{cusp-native-full-flux-number-operator}
\end{equation}
The simultaneous vacuum
\(\ker a_+\cap\ker a_-=H^0((F_T,J_{\eta,T}),\mathscr L_{\eta,T})\)
has complex dimension six by
\cref{cusp-native-full-flux-vacuum-dimensions}.

\begin{theorem}[Complete native joint Landau decomposition]
\label{cusp-native-full-flux-landau-theorem}
For every \(n_+,n_-\in\mathbb Z_{\ge0}\), the joint eigenspace
\begin{equation}
 \mathcal H_{T;n_+,n_-}
 =\ker(N_+-n_+)\cap\ker(N_--n_-)
 \label{cusp-native-full-flux-joint-space}
\end{equation}
has complex dimension six and energy
\begin{equation}
 E_{T;n_+,n_-}
 =(2n_++1)\mathcal B_{T,+}
 +(2n_-+1)\mathcal B_{T,-}.
 \label{cusp-native-full-flux-joint-energy}
\end{equation}
Moreover
\begin{equation}
 L^2(F_T,\mathscr L_{\eta,T})
 =\widehat{\bigoplus}_{n_+,n_-\ge0}\mathcal H_{T;n_+,n_-},
 \label{cusp-native-full-flux-complete-decomposition}
\end{equation}
and the map from the vacuum to the joint level is the unitary
\begin{equation}
 \frac{(a_+^*)^{n_+}(a_-^*)^{n_-}}{\sqrt{n_+!n_-!}}.
 \label{cusp-native-full-flux-ladder-isomorphism}
\end{equation}
For a numerical eigenvalue, the multiplicity is six times the number of
joint pairs in \cref{cusp-native-full-flux-joint-energy} producing that
number.
\end{theorem}

\begin{proof}
The elliptic operator has compact resolvent, and the two commuting number
operators preserve its finite-dimensional eigenspaces; they can therefore be
diagonalized simultaneously.  If \(N_\sigma\psi=\lambda_\sigma\psi\), then
\(\lambda_\sigma\ge0\).  Repeated application of \(a_\sigma\) lowers the
Laplacian energy by \(2\mathcal B_{T,\sigma}>0\), so it terminates.  If
\(a_\sigma^r\psi\ne0\) and \(a_\sigma^{r+1}\psi=0\), then
\[
 0=\|a_\sigma^{r+1}\psi\|^2
  =(\lambda_\sigma-r)\|a_\sigma^r\psi\|^2,
\]
and hence \(\lambda_\sigma=r\in\mathbb Z_{\ge0}\).  Lowering both indices
reaches the simultaneous vacuum.

The commutator identities give, for a vacuum vector \(v\),
\[
 \|(a_\sigma^*)^nv\|^2=n!\|v\|^2,
 \qquad a_\sigma^n(a_\sigma^*)^nv=n!v.
\]
On the \(n\)-th level they also give
\((a_\sigma^*)^na_\sigma^n=n!\).  Operators in the two planes commute, so
\cref{cusp-native-full-flux-ladder-isomorphism} is unitary with inverse
\(a_+^{n_+}a_-^{n_-}/\sqrt{n_+!n_-!}\).  This proves the dimension, energy,
and completeness assertions.  Orthogonal joint eigenspaces with coincident
numerical energies must be added, which proves the final sentence.
\end{proof}

The compact-resolvent lowering argument proves that no additional joint
eigenspaces occur.  Independently, this is exactly the case \(m=2\), \(k=1\)
of
\href{https://www.math.hu-berlin.de/~schueth/quantumtori.pdf}%
{Gordon--Kirwin--Schueth--Webb, Proposition~4.5, article pp.~9--10
(volume pp.~175--176)},
\href{https://doi.org/10.1090/pspum/084/1355}%
{doi:10.1090/pspum/084/1355}.  Their signed tensor-power statement gives
every oscillator tuple multiplicity
\(|k|^mV_{E_\eta}=|k|^2\operatorname{Pf}(E_\eta)=6k^2\), so it also supplies
the external completeness check for every \(k\ne0\).

The exact ground energy and first distinct spacing are therefore
\begin{align}
 E_0(T)&=\mathcal B_{T,+}+\mathcal B_{T,-}
 =\frac{2\pi\Xi_T}{D_T}
 =\frac{14\pi}{T}+\frac{\pi(12B-14H)}{T^2}+O(T^{-3}),
 \label{cusp-native-full-flux-ground-energy}\\
 g_0(T)&=2\min\{\mathcal B_{T,+},\mathcal B_{T,-}\}
 =\frac{2\pi(\Xi_T-|S_T|)}{D_T}
 =\frac{48\pi}{\Xi_T+|S_T|}.
 \label{cusp-native-full-flux-global-gap}
\end{align}
On the deep tail \cref{cusp-native-full-flux-deep-tail-window}, this is
\(g_0(T)=2\mathcal B_{T,-}\), with expansion
\begin{equation}
 g_0(T)=\frac{4\pi}{T}-\frac{4\pi H}{T^2}
 +\frac{4\pi}{T^3}\left(H^2-\frac67m^2\right)+O(T^{-4}).
 \label{cusp-native-full-flux-gap-expansion}
\end{equation}

\begin{corollary}[Heat trace and exact fixed-energy count]
\label{cusp-native-full-flux-heat-count}
For \(t>0\),
\begin{equation}
 \operatorname{Tr}(e^{-t\Delta_{\eta,T}})
 =\frac{3}{2\sinh(t\mathcal B_{T,+})
                 \sinh(t\mathcal B_{T,-})}
 =\frac{D_T}{16\pi^2t^2}\bigl(1+O_t(T^{-2})\bigr).
 \label{cusp-native-full-flux-heat-trace}
\end{equation}
For \(E>0\), let \(N_{\eta,T}(E)\) count complex eigenvalues not exceeding
\(E\), with multiplicity.  Then
\begin{equation}
 N_{\eta,T}(E)=
 \begin{cases}
 0,&E<E_0(T),\\[1mm]
 6\displaystyle\sum_{n_+=0}^{\left\lfloor
 (E-E_0(T))/(2\mathcal B_{T,+})\right\rfloor}
 \left(1+\left\lfloor
 \dfrac{E-E_0(T)-2\mathcal B_{T,+}n_+}
       {2\mathcal B_{T,-}}
 \right\rfloor\right),&E\ge E_0(T).
 \end{cases}
 \label{cusp-native-full-flux-counting-function}
\end{equation}
For fixed \(E>0\), as \(T\to\infty\),
\begin{equation}
 N_{\eta,T}(E)
 =\frac{D_T}{32\pi^2}(E-E_0(T))^2+O_E(T)
 =\frac{E^2D_T}{32\pi^2}+O_E(T).
 \label{cusp-native-full-flux-counting-asymptotic}
\end{equation}
\end{corollary}

\begin{proof}
The heat trace is the product of the two geometric series in
\cref{cusp-native-full-flux-joint-energy}, multiplied by six.  Expanding both
hyperbolic sines and using
\cref{cusp-native-full-flux-frequency-product} proves its asymptotic.
The counting formula enumerates the lattice points
\[
 (n_+,n_-)\in\mathbb R_{\ge0}^2,
 \qquad
 2\mathcal B_{T,+}n_++2\mathcal B_{T,-}n_-
 \le E-E_0(T).
\]
Its triangular area is
\((E-E_0(T))^2/(8\mathcal B_{T,+}\mathcal B_{T,-})\).  The difference
between the lattice count and the area is bounded by a constant times the
perimeter plus one, hence is \(O_E(T)\).  Multiplying by six and using the
exact frequency product gives the first asymptotic; the second follows from
\(E_0(T)=O(T^{-1})\) and \(D_T=O(T^2)\).
\end{proof}

\subsection{The split \texorpdfstring{\(SU(2)\)}{SU(2)} critical connection}

Put
\begin{equation}
 \mathscr E_{\eta,T}=\mathscr L_{\eta,T}\oplus
 \mathscr L_{\eta,T}^{-1},
 \qquad
 A_{\eta,T}=\nabla_{\eta,T}\oplus\nabla_{\eta,T}^{-1}.
 \label{cusp-native-full-flux-split-bundle-connection}
\end{equation}
The determinant is canonically trivial and
\begin{equation}
 c(\mathscr E_{\eta,T})=(1+\eta_T)(1-\eta_T)=1-\eta_T^2,
 \qquad c_2(\mathscr E_{\eta,T})=-\eta_T^2.
 \label{cusp-native-full-flux-Chern-class}
\end{equation}
Thus the orientation-sensitive Chern numbers are
\begin{equation}
 \langle c_2(\mathscr E_{\eta,T}),[F_T]_{\mathrm{mark}}\rangle=-12,
 \qquad
 \langle c_2(\mathscr E_{\eta,T}),
 [F_T]_{x\text{-}\mathrm{complex}}\rangle=+12.
 \label{cusp-native-full-flux-Chern-numbers}
\end{equation}
The cohomology class \(c_2=-\eta_T^2\) is nonzero, so this split
\(SU(2)\)-bundle is not smoothly trivial.  This is another exact distinction
from the unit-flux split bundle of \cref{cusp-magnetic-su2-triviality}.

With
\begin{equation}
 H_{\mathrm{mag}}=\begin{pmatrix}i&0\\0&-i\end{pmatrix},
 \qquad
 \langle X,Y\rangle_{\mathfrak{su}(2)}
 =-\frac12\operatorname{Tr}(XY),
 \label{cusp-native-full-flux-Lie-normalization}
\end{equation}
one has \(\lvert H_{\mathrm{mag}}\rvert=1\) and
\begin{equation}
 F_{A_{\eta,T}}=-2\pi\eta_T H_{\mathrm{mag}}.
 \label{cusp-native-full-flux-SU2-curvature}
\end{equation}
This curvature is parallel; hence
\(d_{A_{\eta,T}}^*F_{A_{\eta,T}}=0\), and \(A_{\eta,T}\) is Yang--Mills.
For
\begin{equation}
 \operatorname{YM}_T(A)=\frac1{2g_{\mathrm{YM}}^2}
 \int_{F_T}|F_A|^2\,\mathrm d\operatorname{vol}_{g_T},
 \label{cusp-native-full-flux-action-normalization}
\end{equation}
diagonalizing \cref{cusp-native-full-flux-K} gives the exact action
\begin{align}
 \operatorname{YM}_T(A_{\eta,T})
 &=\frac{D_T}{2g_{\mathrm{YM}}^2}
   (\mathcal B_{T,+}^2+\mathcal B_{T,-}^2)
 \label{cusp-native-full-flux-action-frequency}\\
 &=\frac{2\pi^2}{g_{\mathrm{YM}}^2D_T}
   (q_T^2+36L_T^2+72m_T^2)
 \label{cusp-native-full-flux-action-coordinate}\\
 &=\frac{24\pi^2}{g_{\mathrm{YM}}^2}
 +\frac{2\pi^2S_T^2}{g_{\mathrm{YM}}^2D_T}.
 \label{cusp-native-full-flux-action-excess}
\end{align}
Indeed
\(q_T^2+36L_T^2+72m_T^2=12D_T+(6L_T-q_T)^2\).
Consequently
\begin{equation}
 \operatorname{YM}_T(A_{\eta,T})
 =\frac{74\pi^2}{g_{\mathrm{YM}}^2}
 +\frac{70\pi^2B}{g_{\mathrm{YM}}^2T}+O(T^{-2}).
 \label{cusp-native-full-flux-action-asymptotic}
\end{equation}

\subsection{The closed Jacobi operator and exact Coulomb reduction}

We next define the operator before computing its spectrum.  Let
\begin{equation}
 \mathscr H_T^1=L^2\Omega^1(F_T;\operatorname{ad}\mathscr E_{\eta,T}),
 \qquad
 \mathscr H_T^0=L^2\Omega^0(F_T;\operatorname{ad}\mathscr E_{\eta,T})
 \label{cusp-native-full-flux-real-Hilbert-spaces}
\end{equation}
be real Hilbert spaces.  Their complexifications carry the root
decomposition used below, and every operator constructed here commutes with
the anti-linear real structure.

Define the maximal distributional codifferential by
\begin{align}
 \operatorname{Dom}D_T^{\max}
 =\{a\in\mathscr H_T^1:\;&\text{there is }u\in\mathscr H_T^0\text{ such that}
 \notag\\[-1mm]
 &\langle a,d_{A_{\eta,T}}\phi\rangle_{L^2}
 =\langle u,\phi\rangle_{L^2}
 \text{ for every smooth }\phi\},
 \label{cusp-native-full-flux-max-codiff-domain}\\
 D_T^{\max}a={}&u,
 \qquad \mathscr S_T:=\ker D_T^{\max}.
 \label{cusp-native-full-flux-Coulomb-slice}
\end{align}
Uniqueness follows from density of the test space.  Passing to the limit in
the defining identity proves that \(D_T^{\max}\) is closed, hence the Coulomb
slice \(\mathscr S_T\) is a closed subspace.

On adjoint-valued scalars, let \(\Delta_{A_{\eta,T}}^{(0)}\) be the
nonnegative self-adjoint operator of the closed form
\begin{equation}
 \mathfrak q_T^{(0)}[\phi,\psi]
 =\langle d_{A_{\eta,T}}\phi,d_{A_{\eta,T}}\psi\rangle_{L^2},
 \qquad
 \operatorname{Dom}\mathfrak q_T^{(0)}=H^1\Omega^0.
 \label{cusp-native-full-flux-scalar-form}
\end{equation}
Elliptic regularity gives
\begin{equation}
 \Delta_{A_{\eta,T}}^{(0)}=d_{A_{\eta,T}}^*d_{A_{\eta,T}},
 \qquad
 \operatorname{Dom}\Delta_{A_{\eta,T}}^{(0)}=H^2\Omega^0.
 \label{cusp-native-full-flux-scalar-domain}
\end{equation}

In a constant orthonormal frame, define the symmetric zeroth-order
endomorphism
\begin{equation}
 (V_Ta)_\nu=-2\sum_{\mu=1}^4
 [F_{A_{\eta,T},\nu\mu},a_\mu].
 \label{cusp-native-full-flux-spin-endomorphism}
\end{equation}
Symmetry follows from invariance of the compact Lie-algebra inner product,
antisymmetry of \(F_{\nu\mu}\), and an exchange of \(\nu,\mu\).  Consider
the gauge-fixed form
\begin{equation}
 \mathfrak j_T^{\mathrm{gf}}[a,b]
 =\langle\nabla_{A_{\eta,T}}a,\nabla_{A_{\eta,T}}b\rangle_{L^2}
  +\langle V_Ta,b\rangle_{L^2},
 \qquad
 \operatorname{Dom}\mathfrak j_T^{\mathrm{gf}}=H^1\Omega^1.
 \label{cusp-native-full-flux-Jacobi-form}
\end{equation}

\begin{proposition}[Closed gauge-fixed realization]
\label{cusp-native-full-flux-closed-Jacobi}
The form \cref{cusp-native-full-flux-Jacobi-form} is densely defined, closed,
symmetric, and bounded below.  Its self-adjoint operator
\(\mathcal J_T^{\mathrm{gf}}\) has compact resolvent, domain \(H^2\Omega^1\),
and expression
\begin{equation}
 (\mathcal J_T^{\mathrm{gf}}a)_\nu
 =-\sum_{\mu=1}^4
   \nabla_{A_{\eta,T},e_\mu}\nabla_{A_{\eta,T},e_\mu}a_\nu
  -2\sum_{\mu=1}^4[F_{A_{\eta,T},\nu\mu},a_\mu].
 \label{cusp-native-full-flux-Jacobi-expression}
\end{equation}
It is exactly the Hessian operator with the background-gauge term, before
the overall factor \(g_{\mathrm{YM}}^{-2}\).
\end{proposition}

\begin{proof}
The first summand in \cref{cusp-native-full-flux-Jacobi-form}, together with
the \(L^2\)-norm, is an equivalent \(H^1\)-norm.  The smooth endomorphism
\(V_T\) is bounded and symmetric.  The bounded perturbation theorem for
closed forms gives closedness and a lower bound, and the representation
theorem gives a self-adjoint operator.  Its principal symbol is
\(|\xi|^2\) times the identity.  Elliptic regularity identifies its maximal
\(L^2\)-domain with \(H^2\Omega^1\); direct substitution supplies the reverse
inclusion.  The compact embedding \(H^1\hookrightarrow L^2\) gives compact
resolvent.

For the final assertion, define
\([a\wedge b]_{\mu\nu}=[a_\mu,b_\nu]-[a_\nu,b_\mu]\).  The exact
two-parameter curvature expansion is
\[
 F_{A_{\eta,T}+sa+tb}=F_{A_{\eta,T}}+s d_Aa+t d_Ab
 +\frac{s^2}{2}[a\wedge a]+st[a\wedge b]
 +\frac{t^2}{2}[b\wedge b].
\]
It yields
\begin{equation}
 \operatorname{YM}_T''\big|_{A_{\eta,T}}(a,b)
 =\frac1{g_{\mathrm{YM}}^2}
 \left(\langle d_Aa,d_Ab\rangle
 +\langle F_A,[a\wedge b]\rangle\right).
 \label{cusp-native-full-flux-raw-Hessian}
\end{equation}
On the flat torus, commuting the covariant derivatives with the retained
curvature convention gives
\begin{align*}
 \langle d_Aa,d_Ab\rangle+\langle d_A^*a,d_A^*b\rangle
 &=\langle\nabla_Aa,\nabla_Ab\rangle
 -\sum_{\nu,\mu}\langle a_\nu,[F_{\nu\mu},b_\mu]\rangle,\\
 \langle F_A,[a\wedge b]\rangle
 &=-\sum_{\nu,\mu}\langle a_\nu,[F_{\nu\mu},b_\mu]\rangle.
\end{align*}
The factor \(1/2\) in the two-form inner product is included in the second
identity.  Adding
\(g_{\mathrm{YM}}^{-2}\langle d_A^*a,d_A^*b\rangle\) to
\cref{cusp-native-full-flux-raw-Hessian} gives exactly
\(g_{\mathrm{YM}}^{-2}\mathfrak j_T^{\mathrm{gf}}[a,b]\), including the
coefficient \(-2\) in \cref{cusp-native-full-flux-spin-endomorphism}.
\end{proof}

\begin{proposition}[Resolvent intertwining and the reducing Coulomb slice]
\label{cusp-native-full-flux-Coulomb-reduction}
On smooth sections,
\begin{equation}
 \mathcal J_T^{\mathrm{gf}}d_{A_{\eta,T}}
 =d_{A_{\eta,T}}\Delta_{A_{\eta,T}}^{(0)},
 \qquad
 d_{A_{\eta,T}}^*\mathcal J_T^{\mathrm{gf}}
 =\Delta_{A_{\eta,T}}^{(0)}d_{A_{\eta,T}}^*.
 \label{cusp-native-full-flux-formal-intertwining}
\end{equation}
For \(z\in\mathbb C\setminus\mathbb R\) and
\(f\in\operatorname{Dom}D_T^{\max}\),
\begin{equation}
 D_T^{\max}(\mathcal J_T^{\mathrm{gf}}-z)^{-1}f
 =(\Delta_{A_{\eta,T}}^{(0)}-z)^{-1}D_T^{\max}f.
 \label{cusp-native-full-flux-resolvent-intertwining}
\end{equation}
Thus \(\mathscr S_T\) reduces \(\mathcal J_T^{\mathrm{gf}}\), and its part
\begin{equation}
 \mathcal J_T=\mathcal J_T^{\mathrm{gf}}\big|_{\mathscr S_T},
 \qquad
 \operatorname{Dom}\mathcal J_T=H^2\Omega^1\cap\mathscr S_T
 \label{cusp-native-full-flux-Coulomb-operator}
\end{equation}
is self-adjoint.  Its closed form is the restriction of
\(\mathfrak j_T^{\mathrm{gf}}\) to \(H^1\Omega^1\cap\mathscr S_T\), and the
Yang--Mills Hessian on this slice is
\begin{equation}
 \operatorname{YM}_T''\big|_{A_{\eta,T}}
 =g_{\mathrm{YM}}^{-2}\mathfrak j_T.
 \label{cusp-native-full-flux-Coulomb-Hessian}
\end{equation}
\end{proposition}

\begin{proof}
For a smooth adjoint-valued scalar \(\phi\), write
\(\nabla_\mu=\nabla_{A_{\eta,T},e_\mu}\).  Retaining both curvature terms
when commuting three derivatives gives
\begin{align*}
 &(\mathcal J_T^{\mathrm{gf}}d_A\phi
       -d_A\Delta_A^{(0)}\phi)_\nu\\
 &\quad=-\nabla_\mu\nabla_\mu\nabla_\nu\phi
       +\nabla_\nu\nabla_\mu\nabla_\mu\phi
       -2[F_{\nu\mu},\nabla_\mu\phi]\\
 &\quad=-[\nabla_\mu F_{\mu\nu},\phi]
       =[(d_A^*F_A)_\nu,\phi]=0.
\end{align*}
This proves the first identity in
\cref{cusp-native-full-flux-formal-intertwining}; formal adjunction gives the
second.

Let \(u=(\mathcal J_T^{\mathrm{gf}}-z)^{-1}f\).  The operator domain gives
\(u\in H^2\), hence \(D_T^{\max}u\in H^1\).  Distributionally applying the
second intertwining identity to
\((\mathcal J_T^{\mathrm{gf}}-z)u=f\) yields
\[
 (\Delta_A^{(0)}-z)D_T^{\max}u=D_T^{\max}f.
\]
The right side is in \(L^2\); scalar elliptic regularity upgrades
\(D_T^{\max}u\) to \(H^2\), and the scalar resolvent proves
\cref{cusp-native-full-flux-resolvent-intertwining}.  Both resolvents of the
self-adjoint gauge-fixed operator preserve \(\mathscr S_T\); adjunction shows
that they also preserve its orthogonal complement.  Therefore the slice is
reducing and its operator part is self-adjoint with the displayed domain.
The reducing projection commutes with the spectral calculus, so the form
domain is exactly \(H^1\Omega^1\cap\mathscr S_T\).  The Hessian identity is
the restriction of the last assertion of
\cref{cusp-native-full-flux-closed-Jacobi}.
\end{proof}

\subsection{Charge two, all polarizations, and the Morse index}

Choose root vectors with
\([H_{\mathrm{mag}},E_\pm]=\pm2iE_\pm\).  Then
\begin{equation}
 (\operatorname{ad}\mathscr E_{\eta,T})_{\mathbb C}
 =\underline{\mathbb C}H_{\mathrm{mag}}
 \oplus\mathscr L_{\eta,T}^{\otimes2}E_+
 \oplus\mathscr L_{\eta,T}^{\otimes(-2)}E_-.
 \label{cusp-native-full-flux-root-decomposition}
\end{equation}
The \(E_+\) root connection has curvature \(-4\pi i\eta_T\).  Replacing
\(E_\eta\) by \(2E_\eta\) gives \(J_{2,T}=J_{\eta,T}\), doubles both
frequencies, and changes the vacuum dimension from six to twenty-four by
\cref{cusp-native-full-flux-vacuum-dimensions}.  The \(E_-\) root has
\(q=-2\), uses the conjugate auxiliary structure
\(J_{-2,T}=-J_{\eta,T}\), and has the same positive frequencies and the same
dimension \(6(-2)^2=24\).  Repeating the complete ladder proof of
\cref{cusp-native-full-flux-landau-theorem} therefore gives, for either root,
the joint charge-two scalar levels
\begin{equation}
 E^{(2)}_{T;p,q}
 =2(2p+1)\mathcal B_{T,+}+2(2q+1)\mathcal B_{T,-},
 \qquad p,q\in\mathbb Z_{\ge0},
 \label{cusp-native-full-flux-charge-two-levels}
\end{equation}
each of complex dimension
\begin{equation}
 |\operatorname{Pf}(2E_\eta)|=24.
 \label{cusp-native-full-flux-charge-two-multiplicity}
\end{equation}
Coincident joint pairs must again be summed.

Choose \(O_T\in SO(2)\) with
\(O_T^{\mathsf T}K_TO_T=\operatorname{diag}(\kappa_{T,+},\kappa_{T,-})\),
and define the constant orthonormal coframe
\begin{equation}
 \binom{\mathrm dX_+}{\mathrm dX_-}
 =O_T^{\mathsf T}\binom{\mathrm dx^1}{\mathrm dx^3},
 \qquad
 \binom{\mathrm dY_+}{\mathrm dY_-}
 =O_T^{\mathsf T}\binom{\mathrm dx^2}{\mathrm dx^4}.
 \label{cusp-native-full-flux-orthogonal-coframe}
\end{equation}
Then
\begin{equation}
 \eta_T=\frac{\mathcal B_{T,+}}{2\pi}
          \mathrm dX_+\wedge\mathrm dY_+
         -\frac{\mathcal B_{T,-}}{2\pi}
          \mathrm dX_-\wedge\mathrm dY_-.
 \label{cusp-native-full-flux-normal-form}
\end{equation}
The odd grouping permutation in
\cref{cusp-native-full-flux-orthogonal-coframe} reverses orientation; it is
used only to expose the two blocks and never to infer a Chern sign.

On the \(E_+\) root line, the four constant circular one-forms diagonalize
the spin endomorphism.  Direct substitution into
\cref{cusp-native-full-flux-spin-endomorphism}, with the charge-two root
normalization retained, gives
\begin{center}
\begin{tabular}{>{$}c<{$} >{$}c<{$} >{$}c<{$}}
\toprule
\text{polarization}&V_T\text{-shift}
 &\mathcal J_T^{\mathrm{gf}}\text{-eigenvalue on }(p,q)\\
\midrule
\mathrm dX_+ + i\mathrm dY_+&-4\mathcal B_{T,+}
 &2[(2p-1)\mathcal B_{T,+}+(2q+1)\mathcal B_{T,-}]\\
\mathrm dX_+ - i\mathrm dY_+&+4\mathcal B_{T,+}
 &2[(2p+3)\mathcal B_{T,+}+(2q+1)\mathcal B_{T,-}]\\
\mathrm dX_- - i\mathrm dY_-&-4\mathcal B_{T,-}
 &2[(2p+1)\mathcal B_{T,+}+(2q-1)\mathcal B_{T,-}]\\
\mathrm dX_- + i\mathrm dY_-&+4\mathcal B_{T,-}
 &2[(2p+1)\mathcal B_{T,+}+(2q+3)\mathcal B_{T,-}]\\
\bottomrule
\end{tabular}
\end{center}
The \(E_-\) table is its anti-linear conjugate.  The Cartan sector is the
ordinary nonnegative coefficientwise Laplacian on one-forms.

\begin{theorem}[Exact negative spectrum, walls, and real Morse index]
\label{cusp-native-full-flux-Morse-theorem}
Put
\begin{equation}
 \mathcal B_{T,\max}=\max\{\mathcal B_{T,+},\mathcal B_{T,-}\},
 \quad
 \mathcal B_{T,\min}=\min\{\mathcal B_{T,+},\mathcal B_{T,-}\},
 \quad
 \rho_T^*=\frac{\mathcal B_{T,\max}}{\mathcal B_{T,\min}}
 =\frac{\Xi_T+|S_T|}{\Xi_T-|S_T|}.
 \label{cusp-native-full-flux-ordered-ratio}
\end{equation}
For every regular \(T\), the negative spectrum of the self-adjoint Coulomb
operator is exactly
\begin{equation}
 \left\{
 \nu_{T,n}=2[(2n+1)\mathcal B_{T,\min}-\mathcal B_{T,\max}]
 =\frac{4\pi}{D_T}[n\Xi_T-(n+1)|S_T|]:
 2n+1<\rho_T^*
 \right\}.
 \label{cusp-native-full-flux-general-negative-spectrum}
\end{equation}
Every indicated band has real multiplicity \(48\), and
\begin{equation}
 \operatorname{ind}_{\mathbb R}(\mathcal J_T)
 =48\left\lceil\frac{\rho_T^*-1}{2}\right\rceil.
 \label{cusp-native-full-flux-general-index}
\end{equation}
At the \(n\)-th wall the strict inequality becomes equality.  Equivalently,
\begin{equation}
 \rho_T^*=2n+1
 \iff n\Xi_T=(n+1)|S_T|
 \iff 24n^2D_T=(2n+1)S_T^2.
 \label{cusp-native-full-flux-wall-equations}
\end{equation}
The corresponding branch is then a zero branch and is not counted in the
Morse index.  If one removes the absolute value by choosing the
\(\mathcal B_{T,+}>\mathcal B_{T,-}\) side, the polynomial equation in
\cref{cusp-native-full-flux-wall-equations} must be accompanied by \(S_T>0\);
otherwise squaring introduces the opposite side as an extraneous solution.

On the sufficiently deep tail of
\cref{cusp-native-full-flux-deep-tail-window}, the negative branches are
exactly \(n=0,1,2\),
\begin{equation}
 \nu_{T,n}=\frac{4\pi}{D_T}[n\Xi_T-(n+1)S_T],
 \qquad n=0,1,2,
 \label{cusp-native-full-flux-deep-negative-branches}
\end{equation}
each of real multiplicity \(48\).  Therefore
\begin{equation}
 \operatorname{ind}_{\mathbb R}(\mathcal J_T)=144,
 \label{cusp-native-full-flux-deep-index}
\end{equation}
and
\begin{align}
 \nu_{T,0}&=-\frac{20\pi}{T}+O(T^{-2}),&
 \nu_{T,1}&=-\frac{12\pi}{T}+O(T^{-2}),\notag\\
 \nu_{T,2}&=-\frac{4\pi}{T}+O(T^{-2}),&
 \nu_{T,3}&=+\frac{4\pi}{T}+O(T^{-2}).
 \label{cusp-native-full-flux-negative-asymptotics}
\end{align}
The Hessian operator for
\cref{cusp-native-full-flux-action-normalization} has eigenvalues
\(\nu_{T,n}/g_{\mathrm{YM}}^2\) and the same index.
\end{theorem}

\begin{proof}
Assume first \(\mathcal B_{T,+}>\mathcal B_{T,-}\).  In the first row of
the polarization table, \(p\ge1\) makes the eigenvalue positive.  At \(p=0\)
it is
\[
 2[-\mathcal B_{T,+}+(2q+1)\mathcal B_{T,-}],
\]
which is negative exactly when \(2q+1<\rho_T^*\).  The second and fourth
rows are positive, and the third has minimum
\(2(\mathcal B_{T,+}-\mathcal B_{T,-})>0\).  When
\(\mathcal B_{T,-}>\mathcal B_{T,+}\), the identical argument with the two
planes exchanged shows that only the \(q=0\) part of the third row can be
negative and gives the same formula with the ordered frequencies.  If the
frequencies agree, both candidate lowest branches are zero and none is
negative.  This proves the exhaustive list for the gauge-fixed operator.

Every negative gauge-fixed eigenform is automatically Coulomb.  Elliptic
regularity makes it smooth; if
\(\mathcal J_T^{\mathrm{gf}}a=\lambda a\) with \(\lambda<0\), then
\cref{cusp-native-full-flux-formal-intertwining} gives
\[
 \Delta_{A_{\eta,T}}^{(0)}(d_{A_{\eta,T}}^*a)
 =\lambda d_{A_{\eta,T}}^*a.
\]
Nonnegativity of the scalar operator forces \(d_A^*a=0\).  Hence the full and
Coulomb negative spectral spaces coincide.

One joint charge-two scalar space contributes complex dimension \(24\) in
the \(E_+\) sector.  Its conjugate contributes complex dimension \(24\) in
the complexification of the real adjoint problem.  A complexified real
eigenspace has complex dimension equal to the original real dimension, so
the real multiplicity is \(24+24=48\), not \(24\) or \(96\).

The count of integers \(n\ge0\) satisfying \(2n+1<\rho_T^*\) is
\(\lceil(\rho_T^*-1)/2\rceil\), including the correct exclusion at an odd
integer wall.  Substitution of the ordered frequency formulas gives
\cref{cusp-native-full-flux-general-negative-spectrum}.  The wall equation
is
\(n\Xi_T=(n+1)|S_T|\); squaring and using
\(\Xi_T^2-S_T^2=24D_T\) gives the final polynomial form, and nonnegativity of
both sides gives the converse.  Finally,
\cref{cusp-native-full-flux-deep-tail-window} leaves precisely
\(n=0,1,2\), and the leading frequency terms give
\(\nu_{T,n}=(8n-20)\pi/T+O(T^{-2})\).
\end{proof}

The lowest branch and the action excess are tied by the exact identities
\begin{equation}
 \nu_{T,0}=-2|\mathcal B_{T,+}-\mathcal B_{T,-}|
 =-\frac{4\pi|S_T|}{D_T},
 \qquad
 \operatorname{YM}_T(A_{\eta,T})-\frac{24\pi^2}{g_{\mathrm{YM}}^2}
 =\frac{D_T}{8g_{\mathrm{YM}}^2}\nu_{T,0}^2.
 \label{cusp-native-full-flux-action-instability}
\end{equation}
